ukazuje, že vysoká flexibilita trhu práce může být udržitelná a~sociálně přijatelná, pokud je v~rámci uceleného \enquote{zlatého trojúhelníku} důsledně propojena s~vysokou příjmovou jistotou a~masivní, státem garantovanou \aclins{APZ}. V~komparativní typologii představuje severský model s~vysokou odborovou organizovaností kolem \SI{67}{\percent} a~pokrytím kolektivními smlouvami přesahujícím \SI{80}{\percent}, v~němž vyjednávání na podnikové úrovni nevede k~rozpadu, nýbrž k~řízené adaptaci odvětvových standardů. Absence zákonem stanovené minimální mzdy zde není projevem neoliberální deregulace, nýbrž výsledkem silné pozice sociálních partnerů, kteří kompenzační minima a~pracovní podmínky definují plošně v~rámci odvětvových smluv. Extrémně nízká ochrana zaměstnání (\textit{EPL}), která usnadňuje a~zlevňuje propouštění, tak dánským firmám sice umožňuje bleskově reagovat na odbytové výkyvy či technologické změny, toto riziko je však plně kompenzováno robustním systémem příjmové bezpečnosti (\textit{A-kasse} s~dlouhodobým náhradovým poměrem až \SI{90}{\percent} pro nejzranitelnější skupiny) a~komplexním, kapitálově silným aktivizačním programem (\textit{learnfare}), na nějž Dánsko vynakládá přes \SI{2}{\percent} \acs{HDP}.~\cite{JdeToPruzne_PS1_FlexicurityDK_2019}~\cite{Denmark_LabourMarket}\par

Pro \acacc{ČR} z~toho plyne zásadní strukturální poučení v~kontextu debat o~tzv.~flexinovele zákoníku práce. České politické iniciativy, inspirované doporučeními Národní ekonomické rady vlády (\acs{NERV}) k~nastartování hospodářského růstu, usilují o~převzetí dánského prvku flexibility -- zejména zavedením výpovědi bez udání důvodu, prodloužením zkušebních dob či zkrácením výpovědních lhůt -- avšak zbývající dva pilíře trojúhelníku, bez nichž se koncepce rozpadá, implementují opatrně. Bez uzavření této aktivační a~příjmové mezery a~bez institucionálního posílení kolektivního vyjednávání povede dánská \textit{flexicurity} na českém trhu práce k~extrémní prekarizaci v~evropském srovnání a~specializaci na levnou nekvalifikovanou práci, která v~podmínkách technologické transformace Průmyslu~4.0 a~digitální ekonomiky silné sociální riziko.~\cite{NERV_NavrhyRustu2025}\par
